\begin{lemma}
\label{editorial-source-lemma-support-finite-presentation-constructible}
Let $R$ be a ring and let $M$ be an $R$-module.
If $M$ is a finitely presented $R$-module, then $\text{Supp}(M)$ is a
closed subset of $\Spec(R)$ whose complement is quasi-compact.
\end{lemma}

\begin{proof}
Retain the original finite presentation
$$
R^{\oplus m}\xrightarrow{A}R^{\oplus n}
\xrightarrow{\pi}M\longrightarrow0,
$$
where $A=(a_{ri})$ is the original $n\times m$ matrix.
Write $e_1,\ldots,e_n$ for the original standard basis
of $R^{\oplus n}$ and $a_i$ for column $i$ of $A$.
For each sorted subset $S=\{i_1<\cdots<i_n\}$ of
$\{1,\ldots,m\}$, set
$$
A_S=(a_{i_1},\ldots,a_{i_n}),\qquad
\Delta_S=\det(A_S),\qquad
I=(\Delta_S:|S|=n).
$$
Thus $I$ is exactly the original ideal of maximal minors.

Fix $\mathfrak p\in\Spec(R)$ and retain
$K=\kappa(\mathfrak p)$.
Applying $R\to R_{\mathfrak p}\to K$ to every matrix
entry gives $A_{\mathfrak p}$ and $A(\mathfrak p)$.
The localized presentation is exact. In detail, the last
map sends $v/s$ to $\pi(v)/s$ and is surjective by lifting
the original numerator. If $\pi(v)/s=0$, there is
$u\notin\mathfrak p$ with $u\pi(v)=0$. Then $uv=Aw$
for some $w\in R^{\oplus m}$, and the full equality
$$
\frac vs=\frac{uv}{us}
=A_{\mathfrak p}\left(\frac w{us}\right)
$$
proves its kernel statement. No injectivity of
$R\to R_{\mathfrak p}$ is assumed.
The fibre quotient comparison is
$$
K^n/\operatorname{im}A(\mathfrak p)
\longrightarrow M\otimes_RK,
\qquad
[(\lambda_1,\ldots,\lambda_n)]
\longmapsto\sum_{r=1}^n\pi(e_r)\otimes\lambda_r.
$$
Its inverse sends $m\otimes\lambda$ to
$[\lambda\overline v]$ for a lift $\pi(v)=m$.
Changing $v$ by $Aw$ changes this vector by
$\lambda A(\mathfrak p)\overline w$, which is in the
matrix image. Additivity and $R$-balancing follow from
the original linear map $\pi$. Both composites are
identities on the displayed generators. Hence
$$
\dim_K(M\otimes_RK)=n-\operatorname{rank}_K A(\mathfrak p).
$$

For $n\geq1$, a minor $\Delta_S\notin\mathfrak p$
gives the explicit right inverse
$$
A_{\mathfrak p}
\left(E_S\frac{\operatorname{adj}(A_S)_{\mathfrak p}}
{\Delta_S/1}\right)=I_n,
$$
where $E_S:R^n\to R^m$ sends its $\ell$-th basis
vector to the $i_\ell$-th original basis vector.
Indeed $AE_S=A_S$ and the adjugate identity gives the
displayed product. Thus $M_{\mathfrak p}=0$.
If $M_{\mathfrak p}=0$, the localized map is surjective;
preimages of the original basis vectors reduce to a
right inverse for $A(\mathfrak p)$, so the fibre is zero.
Conversely, if the fibre is zero, choose
$C\in\operatorname{Mat}_{m\times n}(K)$ with
$A(\mathfrak p)C=I_n$. Cauchy--Binet gives
$$
1=\sum_{\substack{S\subseteq\{1,\ldots,m\}\\|S|=n}}
\overline{\Delta_S}\det(C_{S,\{1,\ldots,n\}}).
$$
At least one original minor therefore lies outside
$\mathfrak p$. This proves, with explicit maps, all
the original equivalences
$$
M_{\mathfrak p}\ne0
\quad\Longleftrightarrow\quad
M\otimes_RK\ne0
\quad\Longleftrightarrow\quad
\operatorname{rank}_K A(\mathfrak p)<n.
$$
In particular $\operatorname{Supp}(M)=V(I)$.

If $n=0$, the target is zero and $M=0$. There is one
$0\times0$ minor indexed by the empty subset, with
determinant $1$, so $I=R$ and both sets are empty.
If $m<n$, there are no maximal minors, so $I=0$.
Every fibre then has dimension at least $n-m>0$, giving
support all of $\Spec(R)$. For the zero ring there are
no primes; every module in the presentation is zero,
and the same support identities hold.

There are only finitely many original maximal minors,
and the complement is the exact finite union
$$
\Spec(R)\setminus\operatorname{Supp}(M)
=\bigcup_{\substack{S\subseteq\{1,\ldots,m\}\\|S|=n}}D(\Delta_S).
$$
Each $D(f)$ is quasi-compact. To recall the proof,
localization identifies $\Spec(R_f)$ with $D(f)$:
contraction is one map and
$\mathfrak p\mapsto\{r/f^a:r\in\mathfrak p,\ a\geq0\}$
is its inverse. The prime localization correspondence
gives these inverse maps, and $D(r/f^a)$ corresponds
to $D(r)\cap D(f)$, proving that they are homeomorphisms.
For any ring $B$, refine a cover of $\Spec(B)$ by
basic opens $D(b_\lambda)$. If their ideal were proper,
a maximal ideal containing it would be missed. Hence
$1=\sum_{j=1}^t c_jb_{\lambda_j}$ for finitely many
original cover elements, and those basic opens cover.
This proves quasi-compactness of every spectrum and
of $D(f)$. A finite union of quasi-compact subspaces
is quasi-compact, by restricting any cover to each
of the finitely many subspaces and joining their finite
subcovers. This proves the original assertion. The
empty union is allowed; for $n=0$ the union is $D(1)$.

\medskip\noindent
The presentation also proves the more precise ideal
bounds. Put $H=\operatorname{Ann}_R(M)$. For $n\geq1$,
$$
I\subseteq H,\qquad H^n\subseteq I.
$$
To prove the first, retain the original column indices
and every cofactor sign in the adjugate identity:
$$
\Delta_Se_r
=\sum_{\ell=1}^n(-1)^{r+\ell}
\det A_{\{1,\ldots,n\}\setminus\{r\},\,
S\setminus\{i_\ell\}}\ a_{i_\ell}.
$$
The remaining rows and columns have inherited increasing
order. In coordinate $k$ the right side is the cofactor
expansion along row $r$ of the matrix with that row
replaced by row $k$. It is $\Delta_S$ if $k=r$ and
zero if $k\ne r$, since then two rows coincide.
Applying $\pi$ proves $\Delta_S\pi(e_r)=0$ for all
$r$. These original elements generate $M$, so
$\Delta_S\in H$ and $I\subseteq H$.

For the second bound take arbitrary original
$t_1,\ldots,t_n\in H$. Since $\pi(t_je_j)=0$,
choose $b_j\in R^m$ with $Ab_j=t_je_j$.
Let $B=(b_{ij})$ be the $m\times n$ matrix with these
columns. Then $AB=\operatorname{diag}(t_1,\ldots,t_n)$.
The full ordered-column determinant expansion is
$$
t_1\cdots t_n=\det(AB)
=\sum_{u_1=1}^m\cdots\sum_{u_n=1}^m
\left(\prod_{j=1}^n b_{u_j,j}\right)
\det(a_{u_1},\ldots,a_{u_n}).
$$
Terms with repeated indices vanish because the determinant
has repeated columns. For distinct indices let
$S=\{i_1<\cdots<i_n\}$ be their sorted set and let
$u_j=i_{\sigma(j)}$. The determinant is
$\operatorname{sgn}(\sigma)\Delta_S$. Thus the full sum is
$$
\sum_{\substack{S=\{i_1<\cdots<i_n\}\\S\subseteq\{1,\ldots,m\}}}
\Delta_S
\left(\sum_{\sigma\in\mathfrak S_n}
\operatorname{sgn}(\sigma)\prod_{j=1}^n b_{i_{\sigma(j)},j}\right).
$$
For $\tau=\sigma^{-1}$ the inner product is
$\prod_{\ell=1}^n b_{i_\ell,\tau(\ell)}$, with the same
sign, so the exact Cauchy--Binet identity is
$$
t_1\cdots t_n
=\sum_{\substack{S\subseteq\{1,\ldots,m\}\\|S|=n}}
\Delta_S\det(B_{S,\{1,\ldots,n\}}).
$$
This proves that every product of $n$ possibly different
elements of $H$ is in $I$. Those products generate $H^n$,
so it proves the ideal-power statement, not only the
statement about individual $n$-th powers. If $m<n$
the sum is empty, giving $H^n=0$. If $n=0$, both $H$
and $I$ are $R$ and the convention $H^0=R$ retains the
two containments. They also hold over the zero ring.
For $n\geq1$ they imply $\sqrt I=\sqrt H$: one inclusion
uses $I\subseteq H$; for $h\in H$, the other uses
$h^n\in H^n\subseteq I$, followed by radicals.
Thus the exact support formula is also
$V(I)=V(\operatorname{Ann}_R(M))$.

The exponent $n$ is sharp uniformly. For a field $k$
take the original presentation $A=xI_n$ over $R=k[x]$,
with $m=n\geq1$. It gives
$$
M=(R/(x))^n,\qquad H=(x),\qquad I=(x^n).
$$
Here an element kills all basis quotient classes exactly
when it belongs to $(x)$, and the single maximal minor
is $x^n$. Thus $H^n=I$; for $n\geq2$ the polynomial
$x^{n-1}$ belongs to $H^{n-1}$ and not to $(x^n)$.
Also $I\subsetneq H$ for these $n$.

\medskip\noindent
There is an exact weaker condition for the topological
conclusion. For any original module $N$ and specified
ideal $L$ with $\operatorname{Supp}(N)=V(L)$, its
complement is quasi-compact if and only if
$$
\sqrt L=\sqrt{(f_1,\ldots,f_t)}
\quad\hbox{for finitely many }f_1,\ldots,f_t\in L.
$$
Indeed $D(L)=\bigcup_{f\in L}D(f)$, so quasi-compactness
gives a finite subcover of this specific cover. Taking
complements gives $V(L)=V((f_1,\ldots,f_t))$ and hence
the radical equality. Conversely that equality gives the
same finite union of quasi-compact basic opens.
The empty list is allowed, with generated ideal zero.
The radical characterization used here follows directly:
every prime containing $J$ contains $\sqrt J$. If
$a\notin\sqrt J$, the ring $(R/J)_a$ is nonzero, since
otherwise some original $a^N$ would be in $J$.
A maximal ideal of this localization contracts to a prime
containing $J$ and avoiding $a$. Consequently $\sqrt J$
is exactly the intersection of the primes containing $J$,
which proves the asserted equivalence of radicals and
closed sets without a finiteness assumption on $J$.

Finite presentation is not necessary for this topology.
Take the specified example
$$
R=k[x,y_1,y_2,\ldots],\qquad
I=(x^2,xy_i:i\geq1),\qquad M=R/I,
$$
where $k$ is a field. Here $I\subseteq(x)$ and $x^2\in I$,
while $R/(x)=k[y_1,y_2,\ldots]$ is a domain: two
nonzero polynomials belong to a polynomial ring in
finitely many variables over $k$, where their product
is nonzero. Thus $\sqrt I=(x)$.
The original cyclic quotient has support $V(I)$.
If $a\in I\setminus\mathfrak p$, the unit $a/1$
makes $(R/I)_{\mathfrak p}$ zero; if $I\subseteq\mathfrak p$,
this quotient maps onto the nonzero residue field.
Therefore the support is $V(x)$ and its complement
is $D(x)=D(x^2)$, a quasi-compact basic open. In the
weaker criterion one can take the actual original
element $f_1=x^2\in I$; there is no replacement of
the original ideal by its radical.

To prove that $I$ is not finitely generated, suppose
$I=(h_1,\ldots,h_r)$ and choose $N$ containing all
indices of variables $y_i$ occurring in any $h_j$.
Define $\rho:R\to k[x,z]$ by fixing $x$, sending
$y_{N+1}$ to $z$, and sending every other $y_i$ to
zero. Define also $\rho_0:R\to k[x]$ by fixing $x$
and sending all the $y_i$ to zero. Since $y_{N+1}$
does not occur in $h_j$, one has
$\rho(h_j)=\rho_0(h_j)\in(x^2)$: the last membership
follows by applying $\rho_0$ to the original generators
$x^2,xy_i$ of $I$. Every combination of the proposed
finite generators therefore maps into $(x^2)$.
But $xy_{N+1}\in I$ maps to $xz\notin(x^2)$.
Indeed its coefficient of $x^1$ over $k[z]$ is the
nonzero polynomial $z$, whereas every multiple of
$x^2$ has that coefficient zero. This is a contradiction.

Finally, for every ring and ideal, the cyclic module
$R/I$ is finitely presented exactly when $I$ is
finitely generated. A finite generating list gives
the presentation $R^t\to R\to R/I\to0$ with that
list as the first matrix. Conversely choose a finite
presentation with surjection $\pi:R^b\to R/I$ and
finitely generated kernel $K=(k_1,\ldots,k_t)$.
Let $q:R\to R/I$ be the original quotient. Choose
$v\in R^b$ with $\pi(v)=1+I$, and an $R$-linear
lift $f:R^b\to R$ of $\pi$ by choosing representatives
for the images of the original basis vectors.
Thus $qf=\pi$, $f(v)-1\in I$, and $f(k_j)\in I$.
For each original $a\in I$, the vector $av$ lies
in $K$, so $av=\sum_j c_jk_j$ for some coefficients.
Retaining both original terms gives the exact identity
$$
a=f(av)+a(1-f(v))
=\sum_{j=1}^t c_jf(k_j)+a(1-f(v)).
$$
Every listed generator is in $I$, and the identity
proves the reverse inclusion in
$$
I=(f(k_1),\ldots,f(k_t),1-f(v)).
$$
This covers $R/I=0$ as well, including $b=0$, where
$1-f(v)=1$ supplies the unit. Applying it to the
specified infinitely generated ideal proves that the
example's original $M$ is not finitely presented,
despite its quasi-compact support complement.
\end{proof}